\documentclass[11pt]{article}

\usepackage[utf8]{inputenc}
\usepackage[T1]{fontenc}
\usepackage[french]{babel}
\usepackage[margin=2cm]{geometry}
\usepackage{booktabs}
\usepackage{longtable}
\usepackage{array}
\usepackage{hyperref}
\usepackage{xcolor}

\hypersetup{
    colorlinks=true,
    urlcolor=blue,
    linkcolor=black
}

\newcolumntype{P}[1]{>{\raggedright\arraybackslash}p{#1}}

\title{Datasets LiDAR publics pour le pré-entraînement général \\
\large Stage LIARA -- Reconnaissance d'activités par LiDAR}
\author{Mathieu Bonneau}
\date{28 août 2026}

\begin{document}

\maketitle

\section*{Contexte}

Ce document liste les datasets LiDAR publics identifiés pour constituer le corpus de
pré-entraînement auto-supervisé (approche V-JEPA) du modèle général sur nuages de points,
avant fine-tuning sur le dataset collecté au laboratoire selon le protocole
\texttt{SG\_protocol\_multi\_v1}. Les datasets sont classés par tier de pertinence.

\footnotesize

\section{Tier 1 --- Indoor / humain-centré avec actions ou pose}
\label{sec:tier1}

Datasets prioritaires : environnement intérieur ou proche, sujets humains, actions ou poses
annotées. Le plus proche du protocole du labo.

\begin{longtable}{P{2.6cm} P{3.0cm} P{4.2cm} P{2.4cm} P{3.0cm}}
\toprule
\textbf{Dataset (année)} & \textbf{Capteur / Environnement} & \textbf{Taille \& labels} & \textbf{Licence / accès} & \textbf{Lien} \\
\midrule
\endfirsthead
\toprule
\textbf{Dataset (année)} & \textbf{Capteur / Environnement} & \textbf{Taille \& labels} & \textbf{Licence / accès} & \textbf{Lien} \\
\midrule
\endhead

MM-Fi (2023) &
Ouster OS1 32ch @10Hz -- Indoor &
320,76k frames, 1080 séquences, 40 sujets, 27 classes (14 AVQ + 13 réhabilitation). Pose 2D/3D + action &
Gratuit académique (formulaire) &
\url{https://arxiv.org/abs/2305.10345} \\
\addlinespace

HuCenLife (2023) &
Ouster OS1 128ch -- Indoor + outdoor &
212 séquences, 3931+2254 frames. Actions riches + interactions humain-humain/objet &
Gratuit (GitHub) &
\url{https://github.com/4DVLab/HuCenLife} \\
\addlinespace

L-CAS 3D Point Cloud People (2017) &
Velodyne VLP-16 -- Indoor (Univ. Lincoln) &
28\,002 frames (\textasciitilde{}49 min). Bbox pedestrian/group, pas d'action &
CC BY-NC-SA 4.0 &
\url{https://lcas.lincoln.ac.uk/wp/research/data-sets-software/l-cas-3d-point-cloud-people-dataset/} \\
\addlinespace

SZTAKI-LGA (2016) &
Velodyne HDL-64E -- Outdoor (activités fines) &
28 sujets. Gait + bend/wave/phone call/check watch. Action + ré-identification &
Gratuit (téléchargement direct) &
\url{http://web.eee.sztaki.hu/i4d/SZTAKI-LGA-DB.html} \\
\addlinespace

JRDB (2019--2021) &
2$\times$Velodyne 16ch + 2$\times$Sick -- Indoor + outdoor (campus) &
64 min annotées, 1,8M cuboïdes 3D, 3500 trajectoires. Bbox 3D + pose/action + groupes sociaux &
Gratuit académique &
\url{https://jrdb.stanford.edu/} \\
\addlinespace

HmPEAR (2024) &
LiDAR 128ch -- Outdoor &
300k+ frames, 10 scènes, 25 sujets, 40 actions quotidiennes. Pose SMPL + action &
Académique (demande via OpenReview) &
\url{https://openreview.net/forum?id=hgRElsBV6v} \\
\addlinespace

STCrowd (2022) &
LiDAR 128ch + caméra -- Outdoor (foule dense) &
219k instances piéton, 10\,891 frames, 84 séquences. Bbox + pose. Fort accent occlusion/densité &
Gratuit académique &
\url{https://arxiv.org/abs/2204.01026} \\

\bottomrule
\end{longtable}

\section{Tier 2 --- Humain/piéton pertinent, outdoor ou actions non domestiques}
\label{sec:tier2}

\begin{longtable}{P{2.6cm} P{3.0cm} P{4.2cm} P{2.4cm} P{3.0cm}}
\toprule
\textbf{Dataset (année)} & \textbf{Capteur / Environnement} & \textbf{Taille \& labels} & \textbf{Licence / accès} & \textbf{Lien} \\
\midrule
\endfirsthead
\toprule
\textbf{Dataset (année)} & \textbf{Capteur / Environnement} & \textbf{Taille \& labels} & \textbf{Licence / accès} & \textbf{Lien} \\
\midrule
\endhead

SLOPER4D (2023) &
Ouster OS1 128ch (casque) -- Outdoor urbain &
100k+ frames LiDAR, 12 sujets, 15 séquences. SMPL + trajectoire globale &
Gratuit (GitHub) &
\url{https://github.com/climbingdaily/SLOPER4D} \\
\addlinespace

LiDARHuman26M / LiDARCap (2022) &
LiDAR longue portée (12--28m) -- Outdoor &
184\,048 frames, 13 sujets, 20 mouvements quotidiens. SMPL complet &
Sur demande (papier) &
\url{https://arxiv.org/abs/2203.14698} \\
\addlinespace

Human-M3 (2023) &
4$\times$Livox MID-100 -- Outdoor multi-personnes &
\textasciitilde{}90k enregistrements, 2--10 personnes simultanées. Pseudo-SMPL &
CC (GitHub) &
\url{https://github.com/soullessrobot/Human-M3-Dataset} \\
\addlinespace

CIMI4D (2023) &
LiDAR scan -- Outdoor (murs d'escalade) &
179\,838 frames, 12 grimpeurs, 13 murs. SMPL + interactions homme-scène &
Gratuit académique &
\url{https://arxiv.org/abs/2303.17948} \\
\addlinespace

HSC4D (2022) &
LiDAR porté sur le corps + IMU -- Indoor + outdoor large échelle &
50k frames LiDAR (20Hz). Pose/trajectoire, pas d'action fine &
Gratuit (GitHub) &
\url{https://github.com/climbingdaily/HSC4D} \\
\addlinespace

FreeMotion (2024) &
Multi-LiDAR (studio panoptique) -- Indoor studio &
40 types de mouvement. SMPL multi-vues &
Via papier LiveHPS &
\url{https://arxiv.org/abs/2402.17171} \\
\addlinespace

Waymo Open Dataset -- Perception v2 (pose) (2023+) &
5 LiDAR (1 moyenne + 4 courte portée) -- Outdoor conduite &
Grande échelle, 14 keypoints/piéton. Pose piéton/cycliste, pas d'action domestique &
Gratuit, licence recherche non commerciale &
\url{https://waymo.com/open/data/perception/} \\

\bottomrule
\end{longtable}

\section{Tier 3 --- Corpus génériques volumineux (pré-entraînement bas-niveau, sans humains étiquetés)}
\label{sec:tier3}

\begin{longtable}{P{2.6cm} P{3.0cm} P{4.2cm} P{2.4cm} P{3.0cm}}
\toprule
\textbf{Dataset (année)} & \textbf{Capteur / Environnement} & \textbf{Taille \& labels} & \textbf{Licence / accès} & \textbf{Lien} \\
\midrule
\endfirsthead
\toprule
\textbf{Dataset (année)} & \textbf{Capteur / Environnement} & \textbf{Taille \& labels} & \textbf{Licence / accès} & \textbf{Lien} \\
\midrule
\endhead

ONCE (2021) &
Non précisé (Huawei) -- Outdoor &
1M scènes LiDAR, 144h de conduite, seulement 15k annotées. Pas de labels humains -- \textbf{corpus conçu pour le SSL} &
Gratuit académique, non commercial &
\url{https://arxiv.org/abs/2106.11037} \\
\addlinespace

SemanticKITTI / KITTI-360 (2019/2021) &
Velodyne HDL-64E -- Outdoor &
43\,552 frames / 320k images + 100k scans. Pas de labels humains fins &
CC BY-NC-SA &
\url{https://arxiv.org/abs/1904.01416} \\
\addlinespace

nuScenes (2019) &
Velodyne HDL-32E -- Outdoor &
34\,149 scans, 1000 scènes. Bbox piéton, pas d'action &
Gratuit, non commercial &
\url{https://www.nuscenes.org/} \\
\addlinespace

Argoverse 2 -- Lidar Dataset (2023) &
2$\times$32ch ($\Rightarrow$64ch@10Hz) -- Outdoor &
20\,000 séquences non labellisées, 6M frames &
CC BY-NC-SA &
\url{https://www.argoverse.org/av2.html} \\
\addlinespace

PandaSet (2020) &
Hesai Pandar64 + PandarGT -- Outdoor &
16\,000 sweeps, 103 scènes. Bbox piéton &
CC BY 4.0 (usage commercial autorisé) &
\url{https://pandaset.org/} \\
\addlinespace

Zenseact Open Dataset -- ZOD (2023) &
VLS128 + 2$\times$VLP16 -- Outdoor &
Multi-année, 14 pays européens. Bbox piéton &
CC BY-SA 4.0 &
\url{https://zod.zenseact.com/} \\
\addlinespace

DAIR-V2X (2022) &
40ch (véhicule) + 300ch (infrastructure) -- Outdoor intersections &
71\,254 frames LiDAR. Bbox, pas de pose &
Gratuit académique &
\url{https://arxiv.org/abs/2204.05575} \\
\addlinespace

ApolloScape (2018) &
Non précisé -- Outdoor &
80k scans LiDAR, 1000km de trajectoires. Pas de labels humains &
Gratuit académique &
\url{https://github.com/ApolloScapeAuto/dataset-api} \\
\addlinespace

Lyft Level 5 (2019) $^{\text{\Large\textcolor{orange}{!}}}$ &
Non précisé -- Outdoor &
29k frames, 97 séquences. Bbox piéton &
Statut de maintenance incertain en 2026 &
via Kaggle / archives \\

\bottomrule
\end{longtable}

\normalsize

\section{Référence méthodologique directement pertinente}

\textbf{AD-L-JEPA} (2025, AAAI 2026) est la première adaptation de l'architecture JEPA à des
données LiDAR automobiles : pré-entraînement auto-supervisé par prédiction d'embeddings BEV
(pas de reconstruction pixel/point), testé sur ONCE et Waymo, code public. C'est la preuve de
concept la plus proche du plan V-JEPA~$\rightarrow$~LiDAR envisagé pour ce stage, à étudier en
détail avant d'adapter le code V-JEPA vidéo.
\begin{itemize}
    \item Papier : \url{https://arxiv.org/abs/2501.04969}
    \item Code : \url{https://github.com/HaoranZhuExplorer/adljepa}
\end{itemize}

À noter également \textbf{DeSPITE} (2025), qui propose un apprentissage contrastif
skeleton--point cloud--IMU--texte sur des données proches de MM-Fi.

\section{Recommandation}

Il n'existe presque aucun dataset LiDAR indoor avec des activités domestiques riches --
MM-Fi (Tier~1) reste le seul vraiment comparable au protocole du laboratoire. Stratégie en
deux temps :

\begin{enumerate}
    \item \textbf{Pré-entraînement domaine-proche} : MM-Fi + HuCenLife + JRDB + HmPEAR +
    STCrowd (\textasciitilde{}1M frames cumulées, indoor/humain) $\rightarrow$ apprentissage
    des priors spatio-temporels les plus pertinents pour le domaine cible.
    \item \textbf{Pré-entraînement volume générique} : ONCE (1M scènes non labellisées,
    conçu pour le SSL) + SemanticKITTI/nuScenes/Argoverse~2/PandaSet pour la robustesse
    géométrique générale du nuage de points, malgré le domaine outdoor éloigné du contexte
    domestique.
\end{enumerate}

\textbf{Point de vigilance} : la résolution des capteurs varie fortement d'un dataset à
l'autre (16 à 300 canaux). Il faudra normaliser explicitement la représentation (résolution
angulaire, portée, fréquence) avant le transfert vers l'Unitree L2, sous peine de reproduire
la chute de performance au changement de capteur déjà documentée dans l'article~III de la
revue de littérature (LiDAR Point Cloud Feature Learning with One-Class SVM).

\end{document}
